\section{Conclusão}
\label{sec:conclusao}

Os quatro objetivos do experimento foram atingidos. A rede treinada
sob as restrições do circuito alcançou \SI{97,15}{\percent} de
acurácia no conjunto de teste do MNIST reduzido, com
\SI{96,44}{\percent} de disparo único correto pelo critério do
comparador, e a tradução para hardware é literal: cada um dos
\num{5451} pesos é um resistor da série E96 em um somador inversor, a
ativação é a saturação do próprio amplificador em trilhos calibrados e
a decisão é um comparador por dígito, sem processamento adicional.

Nas simulações da rede completa, o circuito reproduziu o modelo em
software nas saídas que decidem e acertou, com disparo único, as vinte
janelas da sequência de teste tanto com o amplificador ideal (erro
médio de \vErrIdealMedio~\si{\milli\volt}) quanto com o LM741 (erro
médio de \vErrLmMedio~\si{\milli\volt}, concentrado nas saídas
rejeitadas, sem efeito sobre a decisão). As medições internas
confirmaram a correspondência entre a rede neural e o circuito: os
três regimes da curva de transferência aparecem na distribuição das
ativações, e a perda do terra virtual no nó de soma coincide com o
corte da ReLU, como previsto.

A caracterização do LT1810 mostrou que a velocidade do amplificador
moderno tem pré-requisitos. Sem correções, a rede completa é
inutilizável: a capacitância de entrada com os \SI{100}{\kilo\ohm} de
realimentação produz oscilação sustentada em todos os estágios, e o
deslocamento de polarização de \vVtcOffsetMod~\si{\volt} por estágio,
acumulado pelos cinco estágios, faz um comparador arbitrário disparar
independentemente da entrada. Com as duas correções derivadas em
bancada e aplicadas à rede inteira (capacitor de
\SI{2}{\pico\farad} sobre cada realimentação e resistor de equilíbrio
igual à resistência de Thévenin de cada nó de soma), a rede completa
decidiu corretamente todas as janelas, com latência mediana de
\vModRedeLatMediana~\si{\micro\second} e pior caso de
\vModRedeLatMax~\si{\micro\second}, confirmando a estimativa da
bancada de caminho crítico com diferença inferior a
\SI{2}{\percent}. A compensação teve ainda um efeito colateral útil:
ao eliminar a oscilação, reduziu o custo da simulação de
\vModRedeCpuSemComp~para \vModRedeCpuComp~segundos de CPU.

Na comparação com software, a rede com LM741 decide em dezenas de
microssegundos, uma ordem de grandeza atrás da inferência avulsa em
numpy no computador de referência; a rede com LT1810 compensado decide
em menos de \SI{2}{\micro\second} e supera essa inferência avulsa por
fator próximo de cinco. Em processamento por lote o computador supera
qualquer variante do circuito, e a análise de custo indica que a
implementação analógica só se justifica sob requisitos simultâneos de
latência e consumo.

Como continuidade, propõe-se: análise de Monte Carlo com tolerância de
\SI{1}{\percent} nos resistores e offset real dos amplificadores, para
estimar o rendimento de uma montagem física; poda orientada a custo de
hardware, visando uma montagem discreta de bancada; e a medição do
consumo de potência das variantes, completando o critério de
comparação com a implementação em software.
